
% BEGIN S2
 
\par\medskip\noindent\textbf{2 The declarative system and its properties}\par
 
% BEGIN S2.SS1
 
\par\medskip\noindent\textbf{2.1 Definitions}\par
 
% BEGIN S2.SS1.SSS0.P0.SPx1
 
\par\medskip\noindent\textbf{The type system.}\par
 
% BEGIN S2.SS1.SSS0.P0.SPx1.p1
 
Our declarative type system and specification is a version of Martin-Löf’s type theory [40] with its five mutually defined judgments: typing of types, terms, and contexts and conversion (definitional equality) of types and terms, written $\mathrel{\raisebox{-0.5pt}{$\cong$}}$. They are defined as mutual inductive predicates on “raw” objects. We use $=$ to denote meta-level equality – which coincides $\alpha$-equality, as our representation uses pure de Bruijn indices.\par
 
% END S2.SS1.SSS0.P0.SPx1.p1
 
% BEGIN S2.SS1.SSS0.P0.SPx1.p2
 
The language supports dependent function ($\operatorname{\Pi}$) and pair ($\operatorname{\Sigma}$) types, which we call “negative”, with their $\beta$ and $\eta$-rules. We also include inductive types, with dependent large eliminators: an empty type $\operatorname{\bot}$, natural number type $\operatorname{\mathbb{N}}$, and an identity type $\operatorname{\mathbb{I}d}$. Finally, we have a universe $\operatorname{\mathcal{U}}$, with codes for all type formers except itself. We refer to these types as “positive”. Although our work does not directly involve focusing or logical polarity, where the positive/negative distinction stems from, it will still be an important distinction in the way our conversion algorithms work. We focus on $\operatorname{\mathcal{U}}$, $\operatorname{\Pi}$ and $\operatorname{\mathbb{N}}$ as representative examples on paper and direct interested readers to the code for others.\par
 
% END S2.SS1.SSS0.P0.SPx1.p2
 
% BEGIN S2.SS1.SSS0.P0.SPx1.p3
 
The development relies on parallel substitution, written $t\mathord{\left[\sigma\right]}$. If $u$ is a term, we also denote as $u$ the substitution mapping the last variable in context to $u$ and leaving all others untouched, as used in e.g. $\beta$-reduction: $(\operatorname{\lambda}x:A.t)\leavevmode\nobreak\ u\operatorname{\operatorname%
{\rightsquigarrow}^{\star}}t\mathord{\left[u\right]}$. In Rocq, we use AutoSubst 2 [46, 19] to generate boilerplate for the substitution calculus and automate tedious proofs with a decision procedure for meta-level equality of expressions in this calculus.\par
 
% END S2.SS1.SSS0.P0.SPx1.p3
 
% BEGIN S2.F2

\begin{center}\includegraphics[width=0.98\linewidth,height=0.72\textheight,keepaspectratio]{rules-01.pdf}\end{center}
 
Figure 2:   Canonical, neutral, and normal form predicates (excerpt, see [34]). Way to read the BNF-style grammars: $\operatorname{isTy}$ is the predicate on terms given by the grammar $\operatorname{\mathcal{U}}\mid\operatorname{\Pi}x:t.t\mid\operatorname{\mathbb%
{N}}\mid n$ ($n$ stands for a neutral and $t$ for an arbitrary term).\par
 
% END S2.F2
 
% END S2.SS1.SSS0.P0.SPx1
 
% BEGIN S2.SS1.SSS0.P0.SPx2
 
\par\medskip\noindent\textbf{Normal forms and reduction.}\par
 
% BEGIN S2.SS1.SSS0.P0.SPx2.p1
 
Amongst terms, we particularly focus on (weak-head) normal forms, as this is the class with good injectivity properties, defined in Fig. 2. Canonical forms are terms starting with a constructor: $\operatorname{S}$, $\operatorname{\lambda}$, etc. Neutral forms are “stuck” terms, consisting of a spine of eliminators on top of a variable. Normal forms are either. We also define subclasses for each type: $\operatorname{isFun}t$ holds if $t$ is either a neutral or a $\operatorname{\lambda}$-abstraction, and similarly for $\operatorname{isTy}$ and $\operatorname{isNat}$. The $\operatorname{isTy_{+}}$ predicate isolates positive types: $\operatorname{\mathbb{N}}$, $\operatorname{\mathcal{U}}$ and neutrals.\par
 
% END S2.SS1.SSS0.P0.SPx2.p1
 
% BEGIN S2.F3

\begin{center}\includegraphics[width=0.98\linewidth,height=0.72\textheight,keepaspectratio]{rules-02.pdf}\end{center}
 
Figure 3:   Weak-head reduction (excerpt, see [34]).\par
 
% END S2.F3
 
% BEGIN S2.SS1.SSS0.P0.SPx2.p2
 
Finally, reduction $t\operatorname{\operatorname{\rightsquigarrow}^{\star}}t^{\prime}$ is the reflexive, transitive closure of one-step reduction, defined in Fig. 3. We use only a deterministic weak-head reduction. Reduction is not part of the specification of typing, we present it at this stage only to state normalisation in Sec. 2.2.\par
 
% END S2.SS1.SSS0.P0.SPx2.p2
 
% END S2.SS1.SSS0.P0.SPx2
 
% END S2.SS1
 
% BEGIN S2.SS2
 
\par\medskip\noindent\textbf{2.2 Meta-theory}\par
 
% BEGIN S2.SS2.p1
 
Let us now state the main meta-theoretic properties we consider – we comment about how to prove them at the end of this section and in Sec. 5. In Rocq, these are represented as type-classes: most theorems take instances of the relevant type-classes as parameters. This enables fine-grained bookkeeping of the theorems’ dependency, while keeping it lightweight as instances are automatically filled in. This is inspired by the treatment of function extensionality and univalence in the HoTT library [11].\par
 
% END S2.SS2.p1
 
% BEGIN S2.SS2.SSS0.P0.SPx1
 
\par\medskip\noindent\textbf{Type-level injectivity.}\par
 
% BEGIN Thmtheorem1
 
\par\medskip\noindent\textbf{Property 1 (  Injectivity and no-confusion of type constructors).}\par
 
% BEGIN Thmtheorem1.p1
 
Assume $T$ and $T^{\prime}$ are convertible types in weak-head normal form, i.e. we have that $\operatorname{isTy}T$, $\operatorname{isTy}T^{\prime}$, and $\Gamma\vdash T\mathrel{\raisebox{-0.5pt}{$\cong$}}T^{\prime}$. Then one of the following holds:\par
  
\par\noindent  $\blacksquare$ 
% BEGIN S2.I1.i1.p1
 
$T=\operatorname{\mathbb{N}}=T^{\prime}$ or $T=\operatorname{\mathcal{U}}=T^{\prime}$;\par
 
% END S2.I1.i1.p1
 
 
\par\noindent  $\blacksquare$ 
% BEGIN S2.I1.i2.p1
 
$T=\operatorname{\Pi}x:A.B$, $T^{\prime}=\operatorname{\Pi}x:A^{\prime}.B^{\prime}$, with $\Gamma\vdash A^{\prime}\mathrel{\raisebox{-0.5pt}{$\cong$}}A$ and $\Gamma,x:A^{\prime}\vdash B\mathrel{\raisebox{-0.5pt}{$\cong$}}B^{\prime}$;\par
 
% END S2.I1.i2.p1
 
 
\par\noindent  $\blacksquare$ 
% BEGIN S2.I1.i3.p1
 
$T$, $T^{\prime}$ are both neutral, and $\Gamma\vdash T\mathrel{\raisebox{-0.5pt}{$\cong$}}T^{\prime}\operatorname{%
\mathrel{:}}\operatorname{\mathcal{U}}$.\par
 
% END S2.I1.i3.p1
 
  
Any other case is thus impossible (no-confusion), e.g. we cannot have $\Gamma\vdash\operatorname{\Pi}x:A.B\mathrel{\raisebox{-0.5pt}{$\cong$}}%
\operatorname{\mathbb{N}}$.\par
 
% END Thmtheorem1.p1
 
% END Thmtheorem1
 
% BEGIN S2.SS2.SSS0.P0.SPx1.p1
 
This cannot be proven by mere induction, and indeed it implies a form of non-degeneracy, namely equational consistency: not all types are convertible. Consequences are numerous.\par
 
% END S2.SS2.SSS0.P0.SPx1.p1
 
% BEGIN Thmtheorem2
 
\par\medskip\noindent\textbf{Corollary 2 (  Preservation).}\par
 
% BEGIN Thmtheorem2.p1
 
Assuming Property 1, if $\Gamma\vdash A$ and $A\operatorname{\operatorname{\rightsquigarrow}^{\star}}A^{\prime}$ then $\Gamma\vdash A\mathrel{\raisebox{-0.5pt}{$\cong$}}A^{\prime}$. Similarly, if $\Gamma\vdash t\operatorname{\mathrel{:}}A$ and $t\operatorname{\operatorname{\rightsquigarrow}^{\star}}t^{\prime}$ then $\Gamma\vdash t\mathrel{\raisebox{-0.5pt}{$\cong$}}t^{\prime}\operatorname{%
\mathrel{:}}A$.\par
 
% END Thmtheorem2.p1
 
% END Thmtheorem2
 
% BEGIN Thmtheorem3
 
\par\medskip\noindent\textbf{Corollary 3 (  Types classify normal forms).}\par
 
% BEGIN Thmtheorem3.p1
 
Assuming Property 1, consider $t$ a normal form. If $\Gamma\vdash t\operatorname{\mathrel{:}}\operatorname{\Pi}x:A.B$, then $\operatorname{isFun}t$. If $\Gamma\vdash t\operatorname{\mathrel{:}}\operatorname{\mathbb{N}}$, then $\operatorname{isNat}t$. If $\Gamma\vdash t\operatorname{\mathrel{:}}\operatorname{\mathcal{U}}$, then $\operatorname{isTy}t$. If $\Gamma\vdash t\operatorname{\mathrel{:}}T$ and $T$ is neutral, then $t$ is neutral.\par
 
% END Thmtheorem3.p1
 
% END Thmtheorem3
 
% BEGIN S2.SS2.SSS0.P0.SPx1.p2
 
Preservation relies on the injectivity part, while classification uses no-confusion. Corollary 3 is the key lemma towards progress – the fact that every well-typed term is either a normal form or reduces –, although we do not show this in the formalisation. Thus, Property 1 suffices to establish safety in the progress-and-preservation approach [54].\par
 
% END S2.SS2.SSS0.P0.SPx1.p2
 
% END S2.SS2.SSS0.P0.SPx1
 
% BEGIN S2.SS2.SSS0.P0.SPx2
 
\par\medskip\noindent\textbf{Term-level injectivities.}\par
 
% BEGIN S2.SS2.SSS0.P0.SPx2.p1
 
Statements of injectivity (and no-confusion) for term constructors are relatively close to those for types.\par
 
% END S2.SS2.SSS0.P0.SPx2.p1
 
% BEGIN Thmtheorem4
 
\par\medskip\noindent\textbf{Property 4 (  Injectivity and no-confusion at $\operatorname{\mathbb{N}}$).}\par
 
% BEGIN Thmtheorem4.p1
 
Assume $n$ and $n^{\prime}$ are such that $\operatorname{isNat}n$, $\operatorname{isNat}n^{\prime}$, and $\Gamma\vdash n\mathrel{\raisebox{-0.5pt}{$\cong$}}n^{\prime}\operatorname{%
\mathrel{:}}\operatorname{\mathbb{N}}$. Then one of the following must hold:\par
  
\par\noindent  $\blacksquare$ 
% BEGIN S2.I2.i1.p1
 
$n=0=n^{\prime}$;\par
 
% END S2.I2.i1.p1
 
 
\par\noindent  $\blacksquare$ 
% BEGIN S2.I2.i2.p1
 
$n=\operatorname{S}(t)$, $n^{\prime}=\operatorname{S}(t^{\prime})$, with $\Gamma\vdash t\mathrel{\raisebox{-0.5pt}{$\cong$}}t^{\prime}\operatorname{%
\mathrel{:}}\operatorname{\mathbb{N}}$;\par
 
% END S2.I2.i2.p1
 
 
\par\noindent  $\blacksquare$ 
% BEGIN S2.I2.i3.p1
 
$n$, $n^{\prime}$ are both neutral.\par
 
% END S2.I2.i3.p1
 
  
% END Thmtheorem4.p1
 
% END Thmtheorem4
 
% BEGIN S2.SS2.SSS0.P0.SPx2.p2
 
Maybe surprisingly, the eliminator makes these statements directly provable. Indeed, assuming $\operatorname{S}(t)\mathrel{\raisebox{-0.5pt}{$\cong$}}\operatorname{S}(t^{%
\prime})$, by the computation and congruence rule for the eliminator, we get that $t\mathrel{\raisebox{-0.5pt}{$\cong$}}\operatorname{rec}_{\operatorname{\mathbb%
{N}}}(x.\operatorname{\mathbb{N}},\operatorname{S}(t),0,x.y.x)\mathrel{%
\raisebox{-0.5pt}{$\cong$}}\operatorname{rec}_{\operatorname{\mathbb{N}}}(x.%
\operatorname{\mathbb{N}},\operatorname{S}(t^{\prime}),0,x.y.x)\mathrel{%
\raisebox{-0.5pt}{$\cong$}}t^{\prime}\operatorname{\mathrel{:}}\operatorname{%
\mathbb{N}}$. We can similarly reduce no-confusion to that for types. A similar phenomenon happens at $\operatorname{\Pi}$ types, where injectivity of $\operatorname{\lambda}$ can also be directly derived.\par
 
% END S2.SS2.SSS0.P0.SPx2.p2
 
% BEGIN S2.SS2.SSS0.P0.SPx2.p3
 
Injectivity for the universe ( ) follows a similar pattern, but is not so easily proven for lack of an eliminator. For neutral types, which have no constructors, Corollary 3 is all we need: two (convertible) normal forms at a neutral type must be neutrals.\par
 
% END S2.SS2.SSS0.P0.SPx2.p3
 
% END S2.SS2.SSS0.P0.SPx2
 
% BEGIN S2.SS2.SSS0.P0.SPx3
 
\par\medskip\noindent\textbf{Neutral injectivities.}\par
 
% BEGIN S2.SS2.SSS0.P0.SPx3.p1
 
To complete the injectivity picture, we turn to neutrals, for which injectivity/no-confusion can be expressed in two ways. The first follows a familiar pattern.\par
 
% END S2.SS2.SSS0.P0.SPx3.p1
 
% BEGIN Thmtheorem5
 
\par\medskip\noindent\textbf{Property 5 (  Injectivity of neutral eliminators).}\par
 
% BEGIN Thmtheorem5.p1
 
Assume $n$ and $n^{\prime}$ are neutrals such that $\Gamma\vdash n\mathrel{\raisebox{-0.5pt}{$\cong$}}n^{\prime}\operatorname{%
\mathrel{:}}T$. Then one of the following must hold:\par
  
\par\noindent  $\blacksquare$ 
% BEGIN S2.I3.i1.p1
 
$n=x=n^{\prime}$, with $(x:T^{\prime})\in\Gamma$ and $\Gamma\vdash T^{\prime}\mathrel{\raisebox{-0.5pt}{$\cong$}}T$;\par
 
% END S2.I3.i1.p1
 
 
\par\noindent  $\blacksquare$ 
% BEGIN S2.I3.i2.p1
 
$n=m\ u$, $n^{\prime}=m^{\prime}\ u^{\prime}$ with $\Gamma\vdash m\mathrel{\raisebox{-0.5pt}{$\cong$}}m^{\prime}\operatorname{%
\mathrel{:}}\operatorname{\Pi}x:A.B$, $\Gamma\vdash u\mathrel{\raisebox{-0.5pt}{$\cong$}}u^{\prime}\operatorname{%
\mathrel{:}}A$ and $\Gamma\vdash B\mathord{\left[u\right]}\mathrel{\raisebox{-0.5pt}{$\cong$}}T$;\par
 
% END S2.I3.i2.p1
 
 
\par\noindent  $\blacksquare$ 
% BEGIN S2.I3.i3.p1
 
$n=\operatorname{rec}_{\operatorname{\mathbb{N}}}(m,x.P,t_{0},x.y.t_{%
\operatorname{S}})$, $n^{\prime}=\operatorname{rec}_{\operatorname{\mathbb{N}}}(m^{\prime},x.P^{%
\prime},t^{\prime}_{0},x.y.t^{\prime}_{\operatorname{S}})$, with convertible subterms.\par
 
% END S2.I3.i3.p1
 
  
% END Thmtheorem5.p1
 
% END Thmtheorem5
 
% BEGIN S2.SS2.SSS0.P0.SPx3.p2
 
However, in some cases this is too strong: at negative types our type-directed algorithm eagerly $\eta$-expands, so the above are only really needed at positive types. To state the “minimal” property we require, we take another way, starting with the following definition.\par
 
% END S2.SS2.SSS0.P0.SPx3.p2
 
% BEGIN Thmtheorem6
 
\par\medskip\noindent\textbf{Definition 6 (  Neutral comparison).}\par
 
% BEGIN Thmtheorem6.p1
 
Neutral comparison (see Sec. A.2), written $\Gamma\vdash n\sim n^{\prime}\operatorname{\mathrel{:}}T$, is the least relation that is:\par
  
\par\noindent  $\blacksquare$ 
% BEGIN S2.I4.i1.p1
 
reflexive on variables, i.e. $\Gamma\vdash x\sim x\operatorname{\mathrel{:}}T$ if $(x:T)\in\Gamma$;\par
 
% END S2.I4.i1.p1
 
 
\par\noindent  $\blacksquare$ 
% BEGIN S2.I4.i2.p1
 
closed under the congruence rules for eliminators;\par
 
% END S2.I4.i2.p1
 
 
\par\noindent  $\blacksquare$ 
% BEGIN S2.I4.i3.p1
 
stable under conversion: if $\Gamma\vdash n\sim n^{\prime}\operatorname{\mathrel{:}}T$ and $\Gamma\vdash T\mathrel{\raisebox{-0.5pt}{$\cong$}}T^{\prime}$ then $\Gamma\vdash n\sim n^{\prime}\operatorname{\mathrel{:}}T^{\prime}$.\par
 
% END S2.I4.i3.p1
 
  
% END Thmtheorem6.p1
 
% END Thmtheorem6
 
% BEGIN S2.SS2.SSS0.P0.SPx3.p3
 
We can then rephrase injectivity of neutral constructors as follows.\par
 
% END S2.SS2.SSS0.P0.SPx3.p3
 
% BEGIN Thmtheorem7
 
\par\medskip\noindent\textbf{Property 7 (  Completeness of neutral comparison).}\par
 
% BEGIN Thmtheorem7.p1
 
Given any two neutrals $n$ and $n^{\prime}$ such that $\Gamma\vdash n\mathrel{\raisebox{-0.5pt}{$\cong$}}n^{\prime}\operatorname{%
\mathrel{:}}T$, we have $\Gamma\vdash n\sim n^{\prime}\operatorname{\mathrel{:}}T$.\par
 
% END Thmtheorem7.p1
 
% END Thmtheorem7
 
% BEGIN Thmtheorem8
 
\par\medskip\noindent\textbf{Proposition 8 ( ).}\par
 
% BEGIN Thmtheorem8.p1
 
Property 5 and Property 7 are equivalent.\par
 
% END Thmtheorem8.p1
 
% END Thmtheorem8
 
% BEGIN S2.SS2.SSS0.P0.SPx3.p4
 
But now, we can weaken it to only consider positive types.\par
 
% END S2.SS2.SSS0.P0.SPx3.p4
 
% BEGIN Thmtheorem9
 
\par\medskip\noindent\textbf{Property 9 (  Completeness of neutral comparison at positive types).}\par
 
% BEGIN Thmtheorem9.p1
 
Given any two neutrals $n$ and $n^{\prime}$ such that $\Gamma\vdash n\mathrel{\raisebox{-0.5pt}{$\cong$}}n^{\prime}\operatorname{%
\mathrel{:}}T$, if moreover $\operatorname{isTy_{+}}T$, then we have $\Gamma\vdash n\sim n^{\prime}\operatorname{\mathrel{:}}T$.\par
 
% END Thmtheorem9.p1
 
% END Thmtheorem9
 
% END S2.SS2.SSS0.P0.SPx3
 
% BEGIN S2.SS2.SSS0.P0.SPx4
 
\par\medskip\noindent\textbf{Normalisation.}\par
 
% BEGIN S2.SS2.SSS0.P0.SPx4.p1
 
The last property is normalisation. Its main use is to argue for termination of the algorithm, so it has to go beyond weak-head reduction and take into account $\eta$-expansion, and subterms of (weak-head) normal forms.\par
 
% END S2.SS2.SSS0.P0.SPx4.p1
 
% BEGIN Thmtheorem10
 
\par\medskip\noindent\textbf{Definition 10 (  Deeply normalising).}\par
 
% BEGIN Thmtheorem10.p1
 
A term $t$ is deeply normalising at type $A$ in context $\Gamma$ if they reduce respectively to weak-head normal forms $t^{\prime}$ and $A^{\prime}$, and moreover:\par
  
\par\noindent  $\blacksquare$ 
% BEGIN S2.I5.i1.p1
 
$A^{\prime}$ is a negative type, and the $\eta$-expansion of $t^{\prime}$ is itself normalising;\par
 
% END S2.I5.i1.p1
 
 
\par\noindent  $\blacksquare$ 
% BEGIN S2.I5.i2.p1
 
or $A^{\prime}$ is a positive type, and the subterms of $t^{\prime}$ are normalising at the relevant types.\par
 
% END S2.I5.i2.p1
 
  
A type $A$ is deeply normalising in context $\Gamma$ if it reduces to a weak-head normal form, whose subterms are themselves normalising. See [34] for the inference rules.\par
 
% END Thmtheorem10.p1
 
% END Thmtheorem10
 
% BEGIN S2.SS2.SSS0.P0.SPx4.p2
 
While not apparent in this short version of the definition, the context actually plays a role in order to express the type at which a neutral should be normalising.\par
 
% END S2.SS2.SSS0.P0.SPx4.p2
 
% BEGIN Thmtheorem11
 
\par\medskip\noindent\textbf{Property 11 (  Deep normalisation).}\par
 
% BEGIN Thmtheorem11.p1
 
Every well-typed term is deeply normalising (at its type). Every well-formed type is deeply normalising.\par
 
% END Thmtheorem11.p1
 
% END Thmtheorem11
 
% BEGIN S2.SS2.SSS0.P0.SPx4.p3
 
Meta-theoretic consequences include canonicity ( ) – any closed term of type $\operatorname{\mathbb{N}}$ is convertible to one of the form $\operatorname{S}^{n}0$ – and, most importantly, logical consistency. This means that to prove Property 11, the meta-theory must be logically stronger than the object theory.\par
 
% END S2.SS2.SSS0.P0.SPx4.p3
 
% BEGIN Thmtheorem12
 
\par\medskip\noindent\textbf{Corollary 12 (  Logical consistency).}\par
 
% BEGIN Thmtheorem12.p1
 
Assuming Property 11, there are no closed terms of type $\operatorname{\bot}$.\par
 
% END Thmtheorem12.p1
 
% END Thmtheorem12
 
% END S2.SS2.SSS0.P0.SPx4
 
% BEGIN S2.SS2.SSS0.P0.SPx5
 
\par\medskip\noindent\textbf{Proving the properties.}\par
 
% BEGIN S2.SS2.SSS0.P0.SPx5.p1
 
In our formalisation, most properties are proven via a logical relation, in the style pioneered by Abel et al. [3], and later adapted in Rocq [7], to which we refer for details. Multiple universes in the meta-theory provide the logical strength we need to prove consistency. This approach is however by far not the only possibility, see the discussion in Sec. 5. Note the following two points.\par
 
% END S2.SS2.SSS0.P0.SPx5.p1
 
% BEGIN S2.SS2.SSS0.P0.SPx5.p2
 
First, the proof of deep normalisation is indirect. Indeed, the logical relation is parameterised by an abstract partial equivalence relation (PER) relating neutrals, which we instantiate with algorithmic neutral comparison to obtain normalisation for subterms of neutrals. To avoid this detour through algorithmic aspects, we envisioned strengthening the logical relation to directly encompass neutrals, but did not find a satisfying approach.\par
 
% END S2.SS2.SSS0.P0.SPx5.p2
 
% BEGIN S2.SS2.SSS0.P0.SPx5.p3
 
Second, as logical relations at negative types are naturally defined by observations (e.g. a function is reducible when its application to reducible inputs are reducible), these types “know nothing” about neutrals, and completeness of neutral comparison at negative types cannot be derived easily. We expect other proofs of “semantic” flavour to have similar issues with full completeness of neutrals, making the restriction to positive types (Property 9 vs. Property 7) an important distinction. And indeed, while completeness at positive types is a direct consequence of our logical relation ( ), the approach we found to completeness at all types takes a long detour via the two algorithmic presentations ( ) and strengthening.\par
 
% END S2.SS2.SSS0.P0.SPx5.p3
 
% END S2.SS2.SSS0.P0.SPx5
 
% END S2.SS2
 
% END S2


% BEGIN S3
 
\par\medskip\noindent\textbf{3 Algorithmic typing and conversion}\par
 
% BEGIN S3.p1
 
Let us now turn to our algorithms for conversion and typing. Only the untyped conversion algorithm is new, the typed conversion and bidirectional typing ones come directly from MLTT à la Coq [7]. We give the algorithms informally as rules, which are translated in the formalisation in two ways. First ( ) as inductively defined relations, like declarative typing. Second, we implement them as functions ( ). As we wish to separate their definition from their termination argument, we must embrace partiality. To that end, we use McBride’s free recursion monad [41], via Winterhalter’s PartialFun [53]. This lets us define the functions prior to any proofs; immediately execute them with bounded recursion depth – McBride’s “petrol semantics”; extract a typing (dis)proof from a terminating execution once we prove either soundness; and get proper decidability once we show normalisation. Cherry on the cake: as the recursion monad exposes the programs’ recursive structure, we can generically derive “functional” induction principles, following said recursive structures.\par
 
% END S3.p1
 
% BEGIN S3.SS0.SSS0.P0.SPx1
 
\par\medskip\noindent\textbf{Bidirectional typing.}\par
 
% BEGIN S3.SS0.SSS0.P0.SPx1.p1
 
Our typing algorithm, as most for dependent types, is bidirectional: it separates inference, written $\Gamma\vdash t\mathrel{\rhd}T$, where the type is an output to be found, from type checking, written $\Gamma\vdash t\operatorname{\mathrel{\lhd}}T$, where the type is a constraint given as input. We refer to others [42, 20, 31] for exposition – this last reference is the closest to us. Bidirectional typing generally has two main interests, although only the first is relevant to us here.\par
 
% END S3.SS0.SSS0.P0.SPx1.p1
 
% BEGIN S3.SS0.SSS0.P0.SPx1.p2
 
First, bidirectional typing gives better control over conversion. By replacing the conversion rule with syntax-directed rules, we get a term-directed system, amenable to implementation. This makes the bidirectional system a priori more restrictive than the declarative one, since we have strictly constrained the way conversion can be used. The main point of negative soundness is to show that this restriction does not, in fact, lose anything.\par
 
% END S3.SS0.SSS0.P0.SPx1.p2
 
% BEGIN S3.SS0.SSS0.P0.SPx1.p3
 
Second, propagating types “up” in derivations lightens the need for type annotations and gives better error messages. This is important for surface languages, less for kernels, where terms are typically not human-written. On the contrary, extra information can be useful in various ways, and we would rather dispense of the complication of segregating terms containing too little information to infer. Hence, as Rocq and Lean, we adopt a syntax where all terms infer, although completeness for the other approach is certainly interesting and feasible [35].\par
 
% END S3.SS0.SSS0.P0.SPx1.p3
 
% BEGIN S3.F4

\begin{center}\includegraphics[width=0.98\linewidth,height=0.72\textheight,keepaspectratio]{rules-03.pdf}\end{center}
 
Figure 4:   Typed conversion algorithm (excerpt, see Sec. B.1).\par
 
% END S3.F4
 
% END S3.SS0.SSS0.P0.SPx1
 
% BEGIN S3.SS0.SSS0.P0.SPx2
 
\par\medskip\noindent\textbf{Typed conversion is bidirectional.}\par
 
% BEGIN S3.SS0.SSS0.P0.SPx2.p1
 
A point already implicitly present in Abel et al. [3], but made explicit in Adjedj et al. [7], is that the type-directed conversion algorithm is bidirectional, too. It is decomposed in two main functions, one to compare general terms and the other dedicated to neutrals. The former takes a type as input – it is checking –, while the latter reconstructs types as it traverses the neutrals – it is inferring.\par
 
% END S3.SS0.SSS0.P0.SPx2.p1
 
% BEGIN S3.SS0.SSS0.P0.SPx2.p2
 
More precisely, conversion (see Fig. 4), written $\Gamma\vdash t\mathrel{\raisebox{-0.5pt}{$\cong$}}u\operatorname{\mathrel{\lhd%
}}A$, goes as follows:\par
  
\par\noindent  1. 
% BEGIN S3.I1.i1.p1
 
reduce $t$, $u$ and $A$ to weak-head normal forms;\par
 
% END S3.I1.i1.p1
 
 
\par\noindent  2. 
% BEGIN S3.I1.i2.p1
 
if $A$ is a negative type ($\operatorname{\Pi},\operatorname{\Sigma}$), recursively compare the $\eta$-expansions of $t$ and $u$;\par
 
% END S3.I1.i2.p1
 
 
\par\noindent  3. 
% BEGIN S3.I1.i3.p1
 
if $A$ is a positive type ($\operatorname{\mathbb{N}},\operatorname{\mathcal{U}},\dots$) and the terms are canonical, use the relevant congruence;\par
 
% END S3.I1.i3.p1
 
 
\par\noindent  4. 
% BEGIN S3.I1.i4.p1
 
if $A$ is a positive type and the terms are neutrals, call neutral comparison.\par
 
% END S3.I1.i4.p1
 
  
Steps 2-4 are delegated to an auxiliary function, denoted $\mathrel{\raisebox{-0.5pt}{$\cong$}}_{\mathrm{h}}$ (the “h” is for “head”). Neutral comparison, written $\Gamma\vdash n\sim n^{\prime}\mathrel{\rhd}S$ follows their structure, applying congruences as it goes.\par
 
% END S3.SS0.SSS0.P0.SPx2.p2
 
% BEGIN S3.SS0.SSS0.P0.SPx2.p3
 
Rule   Typed conversion algorithm (excerpt, see Sec.˜B.1). demonstrates this bidirectional yoga. As conversion is type-directed, we need a type $A$ to compare the arguments, that is obtained through the recursive comparison of the functions, which must thus be inferring. We would not be able to construct $\operatorname{\Pi}x:A.B$ anyway, even given $B\mathord{\left[u\right]}$: inverting substitutions is unfeasible. Thus, even with a syntax with complete inference, the distinction between checking constructors and inferring eliminators, a landmark of the “Nordic” approach to bidirectionalism [43, 42], appears naturally.\par
 
% END S3.SS0.SSS0.P0.SPx2.p3
 
% END S3.SS0.SSS0.P0.SPx2
 
% BEGIN S3.SS0.SSS0.P0.SPx3
 
\par\medskip\noindent\textbf{Untyped conversion.}\par
 
% BEGIN S3.F5

\begin{center}\includegraphics[width=0.98\linewidth,height=0.72\textheight,keepaspectratio]{rules-04.pdf}\end{center}
 
Figure 5:   Untyped algorithmic conversion (excerpt, see Sec. B.2).\par
 
% END S3.F5
 
% BEGIN S3.SS0.SSS0.P0.SPx3.p1
 
Untyped conversion (Fig. 5) looks much simpler than typed conversion: there are no types, no separation between type and term conversion, and the separation of neutral comparison is kept mostly for readability.\par
 
% END S3.SS0.SSS0.P0.SPx3.p1
 
% BEGIN S3.SS0.SSS0.P0.SPx3.p2
 
The main point is that we replace the all-encompassing   Typed conversion algorithm (excerpt, see Sec.˜B.1). by term-directed rules, an approach pioneered by Coquand, first for functions [16], then pairs [2]. When comparing $\operatorname{\lambda}$-abstractions, we apply the congruence   Untyped algorithmic conversion (excerpt, see Sec.˜B.2). – which would also happen in   Typed conversion algorithm (excerpt, see Sec.˜B.1). after contracting the redexes. If one side is a $\operatorname{\lambda}$ and the other a neutral (  Untyped algorithmic conversion (excerpt, see Sec.˜B.2).), we again simulate   Typed conversion algorithm (excerpt, see Sec.˜B.1)., but this time we see the application to a fresh variable is visible on the neutral’s side. And this is all! That is, there is no $\eta$-rule when the two sides are neutrals. How could there be? In the extreme case of two variables, the only information we could use is… type information, which we precisely decided not to rely on. This discrepancy on neutrals at negative types is the key difference between the two algorithms.\par
 
% END S3.SS0.SSS0.P0.SPx3.p2
 
% END S3.SS0.SSS0.P0.SPx3
 
% BEGIN S3.SS0.SSS0.P0.SPx4
 
\par\medskip\noindent\textbf{Comparison.}\par
 
% BEGIN S3.SS0.SSS0.P0.SPx4.p1
 
As a concrete example, we can look at the “execution trace” in the following problem: $x:\operatorname{\mathbb{N}}\to\operatorname{\mathbb{N}}\vdash x\mathrel{%
\raisebox{-0.5pt}{$\cong$}}x\operatorname{\mathrel{:}}\operatorname{\mathbb{N}%
}\to\operatorname{\mathbb{N}}$. The untyped algorithm immediately uses   Untyped algorithmic conversion (excerpt, see Sec.˜B.2). to go in “neutral mode”, and finishes with   Untyped algorithmic conversion (excerpt, see Sec.˜B.2).. The typed algorithm, on the other hand, first $\eta$-expands, introducing a fresh variable $y$, and recursively compares $x\ y$ to $x\ y$. Only then, by reaching a positive type, does it change to “neutral mode”. Neutral comparison then peels off the application node, compares $x$ to $x$ again but this time as neutrals, and finally applies variable reflexivity. It also needs to compare $y$ to $y$ first as normal forms at $\operatorname{\mathbb{N}}$, then as neutrals, and finally terminates. We can already see the work spared by the term-directed algorithm – exacerbated for $\operatorname{\Sigma}$ types, where $\eta$ introduces duplication.\par
 
% END S3.SS0.SSS0.P0.SPx4.p1
 
% END S3.SS0.SSS0.P0.SPx4
 
% END S3


% BEGIN S4.SS1
 
\par\medskip\noindent\textbf{4.1 Properties of typed conversion}\par
 
% BEGIN S4.SS1.p1
 
Most results in this section are not per se new, although they were significantly reworked, in particular to precisely track their dependencies. In particular, termination relies explicitly on normalisation, rather than on the reflexivity of the logical relation. The work around negative soundness is entirely new.\par
 
% END S4.SS1.p1
 
% BEGIN S4.SS1.SSS0.P0.SPx1
 
\par\medskip\noindent\textbf{Positive soundness.}\par
 
% BEGIN S4.SS1.SSS0.P0.SPx1.p1
 
We split the proof in two, using the inductive judgment as intermediate. The first half is almost tautological, and does not depend on any meta-theory.\par
 
% END S4.SS1.SSS0.P0.SPx1.p1
 
% BEGIN Thmtheorem13
 
\par\medskip\noindent\textbf{Proposition 13 (  Positive soundness of the typed conversion function).}\par
 
% BEGIN Thmtheorem13.p1
 
If a typed conversion function returns a positive answer, then the corresponding judgment of Fig. 4 holds.\par
 
% END Thmtheorem13.p1
 
% END Thmtheorem13
 
% BEGIN Thmtheorem14
 
\par\medskip\noindent\textbf{Proposition 14 (  Positive soundness of the typed algorithmic conversion judgment).}\par
 
% BEGIN Thmtheorem14.p1
 
Assuming injectivity of type constructors (Property 1), the algorithmic typed conversion judgments imply their postconditions whenever their preconditions hold.\par
 
% END Thmtheorem14.p1
 
% END Thmtheorem14
 
% BEGIN S4.SS1.SSS0.P0.SPx1.p2
 
This second half is the interesting one. The crucial part of this proof, and the backbone of others in this section, is to show that each inductive rule correctly preserves the invariants. This is a very regular property, so much that we meta-programmed444 We hardcoded these in the current code to remove a dependency which caused proof engineering issues.  the generation of “local preservation lemmas” ( ) for each rule, expressing this idea. We assemble these lemmas into a “bundled induction principle”, again programmatically generated, which threads the invariants through, giving access to them as extra hypotheses in each induction step. Soundness follows by instantiating this induction principle with the constant True predicates.\par
 
% END S4.SS1.SSS0.P0.SPx1.p2
 
% END S4.SS1.SSS0.P0.SPx1
 
% BEGIN S4.SS1.SSS0.P0.SPx2
 
\par\medskip\noindent\textbf{Negative soundness.}\par
 
% BEGIN S4.SS1.SSS0.P0.SPx2.p1
 
For negative soundness, we cannot go through the intermediate algorithmic judgments, as they only encode successful runs of the functions. Thus, we bite the bullet and directly relate the output of the functions to the declarative judgments.\par
 
% END S4.SS1.SSS0.P0.SPx2.p1
 
% BEGIN Thmtheorem15
 
\par\medskip\noindent\textbf{Proposition 15 (  Negative soundness of typed algorithmic conversion).}\par
 
% BEGIN Thmtheorem15.p1
 
Assuming injectivity of type and term constructors (Properties 1, 4, etc.), and completeness of neutral conversion at positive types (Property 9), if a typed conversion function returns a negative answer and its precondition holds, then its postcondition cannot hold.\par
 
% END Thmtheorem15.p1
 
% END Thmtheorem15
 
% BEGIN S4.SS1.SSS0.P0.SPx2.p2
 
Again, we rely on invariant preservation, the proof is otherwise straightforward: we have extracted exactly the injectivity properties needed to justify each step. In essence, we show that at any stage the current “remaining goals” are equivalent to the initial query, relying on injectivity to know that we only use invertible congruences, which are safe to apply.\par
 
% END S4.SS1.SSS0.P0.SPx2.p2
 
% END S4.SS1.SSS0.P0.SPx2
 
% BEGIN S4.SS1.SSS0.P0.SPx3
 
\par\medskip\noindent\textbf{Termination.}\par
 
% BEGIN S4.SS1.SSS0.P0.SPx3.p1
 
As expected, this is where normalisation becomes necessary. As before, we need invariants throughout, and thus injectivity of type constructors. The most subtle part of the proof concerns the reduction machine.\par
 
% END S4.SS1.SSS0.P0.SPx3.p1
 
% BEGIN Thmtheorem16
 
\par\medskip\noindent\textbf{Lemma 16 (  Completeness of reduction).}\par
 
% BEGIN Thmtheorem16.p1
 
Assuming Property 1, if $t$ is well-typed and reduces to the weak-head normal form $t^{\prime}$, then the reduction function applied to $t$ returns $t^{\prime}$.\par
 
% END Thmtheorem16.p1
 
% END Thmtheorem16
 
% BEGIN S4.SS1.SSS0.P0.SPx3.p2
 
The main task is to define a well-founded order on the configurations of the abstract machine implementing reduction, a lexicographic product of subterm ordering and reduction. The assumption that $t$ reduces to a weak-head normal form ensures that this order is well-founded. Winterhalter introduced a similar order for MetaCoq [52, chap. 23].\par
 
% END S4.SS1.SSS0.P0.SPx3.p2
 
% BEGIN Thmtheorem17
 
\par\medskip\noindent\textbf{Proposition 17 (  Termination of typed algorithmic conversion).}\par
 
% BEGIN Thmtheorem17.p1
 
Assuming injectivity of type constructors (Property 1) and deep normalisation (Property 11), typed conversion functions terminate whenever their preconditions are satisfied.\par
 
% END Thmtheorem17.p1
 
% END Thmtheorem17
 
% BEGIN S4.SS1.SSS0.P0.SPx3.p3
 
The proof is by induction on deep normalisation, again using local preservation lemmas. The inductive structure of deep normalisation readily encodes the combination of subterm, reduction and $\eta$-expansion that one needs to be well-founded. In contrast, in MetaCoq, Winterhalter introduced a complex dependent lexicographic order modulo [52, chap 24], and admits that his setup would not easily incorporate $\eta$-expansion.\par
 
% END S4.SS1.SSS0.P0.SPx3.p3
 
% END S4.SS1.SSS0.P0.SPx3
 
% BEGIN S4.SS1.SSS0.P0.SPx4
 
\par\medskip\noindent\textbf{Non-triviality.}\par
 
% BEGIN S4.SS1.SSS0.P0.SPx4.p1
 
Soundness, both positive and negative, has a default: the nowhere-defined function is sound! Although we repeat that these properties should not be read to merely say “there exists a function such that…” but to characterise a particular checker – that we can run on examples ( ) to check for non-triviality –, it is good to formally ensure that our functions are non-trivial. Again, this is split in two.\par
 
% END S4.SS1.SSS0.P0.SPx4.p1
 
% BEGIN Thmtheorem18
 
\par\medskip\noindent\textbf{Proposition 18 (  Completeness of the implementation).}\par
 
% BEGIN Thmtheorem18.p1
 
Assuming injectivity of type constructors (Property 1), if an algorithmic conversion judgment and its precondition hold, then the corresponding function answers positively.\par
 
% END Thmtheorem18.p1
 
% END Thmtheorem18
 
% BEGIN S4.SS1.SSS0.P0.SPx4.p2
 
Next, we show that algorithmic conversion satisfies many closure properties, which Abel et al. call “generic conversion”. These are rather rich: all the proofs on the logical relation side can be done with respect to any generic conversion.\par
 
% END S4.SS1.SSS0.P0.SPx4.p2
 
% BEGIN Thmtheorem19
 
\par\medskip\noindent\textbf{Proposition 19 ( ).}\par
 
% BEGIN Thmtheorem19.p1
 
Algorithmic conversion is a generic conversion, thus for instance symmetric and transitive, stable by weakening and expansion, congruent for constructors.\par
 
% END Thmtheorem19.p1
 
% END Thmtheorem19
 
% BEGIN S4.SS1.SSS0.P0.SPx4.p3
 
Compared to declarative typing, the only missing part is congruences for eliminators, which only hold for neutral conversion, and reflexivity – which would be direct by induction if we had all congruence rules. The point is that eliminator congruences are the dangerous ones that can trigger reductions: applying a normal form to a normal form does not yield a normal form. And indeed a term is deeply normalising exactly when it is reflexive for (typed) algorithmic conversion, explaining why reflexivity cannot be achieved easily.\par
 
% END S4.SS1.SSS0.P0.SPx4.p3
 
% END S4.SS1.SSS0.P0.SPx4
 
% BEGIN S4.SS1.SSS0.P0.SPx5
 
\par\medskip\noindent\textbf{Typing.}\par
 
% BEGIN S4.SS1.SSS0.P0.SPx5.p1
 
All of these propositions extend to bidirectional typing. More precisely, assuming any conversion relation/function between types that satisfies the relevant property, we can derive a similar one for typing. This means that the proofs are modular, in that we can swap out a different implementation of conversion, as long as it satisfies the right properties, which is exactly what we are about to do.\par
 
% END S4.SS1.SSS0.P0.SPx5.p1
 
% END S4.SS1.SSS0.P0.SPx5
 
% END S4.SS1


% BEGIN S4.SS2
 
\par\medskip\noindent\textbf{4.2 Properties of untyped conversion}\par
 
% BEGIN S4.SS2.p1
 
The high-level ideas are very similar to typed conversion. For positive soundness there is really no difference: the equivalents of Propositions 13 ( ) and 14 ( ) hold. Even better, we can reuse most of the local preservation lemmas directly.\par
 
% END S4.SS2.p1
 
% BEGIN S4.SS2.p2
 
Negative soundness is where the subtle difference between the two algorithms shows. Echoing the discussion in Secs. 2.2 and 3, the key point is neutrals at negative types. And indeed, while in Proposition 15 we could make do with the weaker completeness (Property 9), this time we really need the full power of completeness at all types.\par
 
% END S4.SS2.p2
 
% BEGIN Thmtheorem20
 
\par\medskip\noindent\textbf{Proposition 20 (  Negative soundness of untyped algorithmic conversion).}\par
 
% BEGIN Thmtheorem20.p1
 
Assuming injectivity of type and term constructors (Properties 1, 4, etc.), and completeness of neutral conversion at all types (Property 7), if the untyped conversion function returns a negative answer and its preconditions holds, then its postcondition cannot hold.\par
 
% END Thmtheorem20.p1
 
% END Thmtheorem20
 
% BEGIN S4.SS2.p3
 
As for termination, there is a twist again, of course at negative types. The issue is that deep normalisation systematically $\eta$-expands terms, but untyped conversion does not always do this. Worse, it is non-deterministic, in the sense that whether a given term is $\eta$-expanded depends on the term it is compared to. Thus, our termination measure is not exactly fit, and circumventing the issue involves proving a form of strengthening of the algorithm: it gives the same output when irrelevant variables are removed from the context.\par
 
% END S4.SS2.p3
 
% BEGIN Thmtheorem21
 
\par\medskip\noindent\textbf{Proposition 21 (  Termination of untyped algorithmic conversion).}\par
 
% BEGIN Thmtheorem21.p1
 
Assuming Property 1 and Property 11, untyped conversion functions terminate whenever their preconditions hold.\par
 
% END Thmtheorem21.p1
 
% END Thmtheorem21
 
% END S4.SS2


% BEGIN S4.SS3
 
\par\medskip\noindent\textbf{4.3 Equivalence between the algorithms}\par
 
% BEGIN S4.SS3.p1
 
To wrap up, we can directly compare the two algorithms. Given how close they are, this should be rather easy, or at any rate easier than comparing declarative systems – where the seemingly innocuous transitivity is actually a source of headaches [44], because it means “real” untyped conversion can relate well-typed terms through ill-typed ones, which cannot be easily emulated by its typed counterpart. In fact, streamlining the relation between typed and untyped conversion was one of our original motivation [31, chap. 6] [32].\par
 
% END S4.SS3.p1
 
% BEGIN Thmtheorem22
 
\par\medskip\noindent\textbf{Proposition 22 (  Typed to untyped conversion).}\par
 
% BEGIN Thmtheorem22.p1
 
Assuming Property 1, the typed algorithmic judgments imply their untyped counterparts whenever their preconditions hold.\par
 
% END Thmtheorem22.p1
 
% END Thmtheorem22
 
% BEGIN S4.SS3.SSS0.P0.SPx1
 
\par\medskip\noindent\textbf{Proof.}\par
 
% BEGIN S4.SS3.SSS0.P0.SPx1.p1
 
The proof is by bundled induction on the typed algorithmic judgment. We simultaneously prove that if $\Gamma\vdash n\mathrel{\raisebox{-0.5pt}{$\cong$}}n^{\prime}\operatorname{%
\mathrel{\lhd}}T$ and moreover $n$, $n^{\prime}$ are neutrals, then $n\sim n^{\prime}$.\par
 
% END S4.SS3.SSS0.P0.SPx1.p1
 
% BEGIN S4.SS3.SSS0.P0.SPx1.p2
 
The main interesting case is of course that of   Typed conversion algorithm (excerpt, see Sec.˜B.1).. In this case, the typing invariants imply, by Corollary 3, that both sides are either abstractions or neutrals. In the first three cases, we can directly conclude by induction hypothesis, up to $\beta$-reduction of the $\eta$-expanded abstraction. Remains the case of two neutrals $n$ and $n^{\prime}$. The extra induction hypothesis for neutrals tells us that $n\ y\sim n^{\prime}\ y$, or more precisely in de Bruijn syntax, that $n\mathord{\left[\uparrow\right]}\ y\sim n^{\prime}\mathord{\left[\uparrow%
\right]}\ y$, where $\uparrow$ is the substitution shifting all indices by 1. We can easily invert this to $n\mathord{\left[\uparrow\right]}\sim n^{\prime}\mathord{\left[\uparrow\right]}$, and conclude by an auxiliary lemma that untyped conversion admits strengthening. $\hfill\blacktriangleleft$ Note the final use of strengthening: while this is readily proven by induction for algorithmic judgments, it is generally difficult to prove for declarative systems! Admitting easy proofs of strengthening is a major advantage of algorithmic presentations [30].\par
 
% END S4.SS3.SSS0.P0.SPx1.p2
 
% BEGIN S4.SS3.SSS0.P0.SPx1.p3
 
For the converse, there is a catch: normalisation is required. Indeed, imagine a type $A$, well-formed in $\Gamma$, such that there is an infinite reduction sequence $A\operatorname{\rightsquigarrow}A_{1}\operatorname{\rightsquigarrow}\dots$ Trying to deduce $\Gamma,x:A\vdash x\mathrel{\raisebox{-0.5pt}{$\cong$}}x\operatorname{\mathrel{%
:}}A$, the untyped algorithm would terminate immediately, while the typed one would diverge on the evaluation of $A$. More insidiously, if a type $B$ keeps on exposing $\operatorname{\Pi}$ constructors – say a solution to $B\mathrel{\raisebox{-0.5pt}{$\cong$}}\operatorname{\mathbb{N}}\to B$ – comparing $x:B$ to itself would terminate immediately on one side, and spin in never ending $\eta$-expansions on the other. We can, however, prove the following crucial lemma.\par
 
% END S4.SS3.SSS0.P0.SPx1.p3
 
% BEGIN Thmtheorem23
 
\par\medskip\noindent\textbf{Lemma 23 (  Completeness of typed neutral conversion at all types).}\par
 
% BEGIN Thmtheorem23.p1
 
Assuming Property 11, given a context $\Gamma$ and $A$, $m$, $n$ respectively a well-formed type and well-typed terms in $\Gamma$, such that $\Gamma\vdash m\sim n\mathrel{\rhd}A^{\prime}$ and $\Gamma\vdash A^{\prime}\mathrel{\raisebox{-0.5pt}{$\cong$}}A$, we have $\Gamma\vdash m\mathrel{\raisebox{-0.5pt}{$\cong$}}n\operatorname{\mathrel{\lhd%
}}{A}$.\par
 
% END Thmtheorem23.p1
 
% END Thmtheorem23
 
% END S4.SS3.SSS0.P0.SPx1
 
% BEGIN S4.SS3.SSS0.P0.SPx2
 
\par\medskip\noindent\textbf{Proof sketch555In the formalisation, we shortcut this proof using the equivalence of declarative and algorithmic typing..}\par
 
% BEGIN S4.SS3.SSS0.P0.SPx2.p1
 
By induction on the deep normalisation of $A$, using which, after a series of $\eta$-expansions of $m$ and $n$, we end up at a positive type. At this point we go into neutral mode, peel off all introduced $\eta$-expansions, and use our neutral comparison hypothesis, adequately weakened to account for the extra introduced variables, to conclude. $\hfill\blacktriangleleft$\par
 
% END S4.SS3.SSS0.P0.SPx2.p1
 
% BEGIN Thmtheorem24
 
\par\medskip\noindent\textbf{Proposition 24 (  Untyped to typed conversion).}\par
 
% BEGIN Thmtheorem24.p1
 
Assuming injectivity of typed constructors and deep normalisation, the untyped algorithmic judgments imply their typed counterparts whenever their preconditions hold.\par
 
% END Thmtheorem24.p1
 
% END Thmtheorem24
 
% BEGIN S4.SS3.SSS0.P0.SPx2.p2
 
After Lemma 23, remains a straightforward induction – as usual, using preservation lemmas as necessary. However, let us emphasise that while the proof of equivalence is relatively direct, it nonetheless relies on difficult properties: strengthening, and normalisation.\par
 
% END S4.SS3.SSS0.P0.SPx2.p2
 
% END S4.SS3.SSS0.P0.SPx2
 
% END S4.SS3


% BEGIN A1
 
\par\medskip\noindent\textbf{Appendix A Declarative Martin-Löf Type Theory}\par
 
% BEGIN A1.p1
 
On top of $t\mathord{\left[u\right]}$ for the substitution of the last variable of $t$ by $u$, we more generally use $t\mathord{\left[u,v,w\right]}$ for the parallel substitution of the last variables of $t$ with $u$, $v$ and $w$.\par
 
% END A1.p1
 
% BEGIN A1.SS1
 
\par\medskip\noindent\textbf{A.1 Declarative conversion}\par
 
% BEGIN A1.SS1.p1
 
\begin{center}\includegraphics[width=0.98\linewidth,height=0.72\textheight,keepaspectratio]{rules-05.pdf}\end{center}
 
% END A1.SS1.p1
 
% END A1.SS1
 
% BEGIN A1.SS2
 
\par\medskip\noindent\textbf{A.2 Neutral comparison}\par
 
% BEGIN A1.SS2.p1
 
\begin{center}\includegraphics[width=0.98\linewidth,height=0.72\textheight,keepaspectratio]{rules-06.pdf}\end{center}
 
% END A1.SS2.p1
 
% END A1.SS2
 
% END A1


% BEGIN A2
 
\par\medskip\noindent\textbf{Appendix B Algorithmic Martin-Löf Type Theory}\par
 
% BEGIN A2.SS1
 
\par\medskip\noindent\textbf{B.1 Typed algorithmic conversion}\par
 
% BEGIN A2.SS1.p1
 
\begin{center}\includegraphics[width=0.98\linewidth,height=0.72\textheight,keepaspectratio]{rules-07.pdf}\end{center}
 
% END A2.SS1.p1
 
% BEGIN A2.SS1.p2
 
\begin{center}\includegraphics[width=0.98\linewidth,height=0.72\textheight,keepaspectratio]{rules-08.pdf}\end{center}
 
% END A2.SS1.p2
 
% BEGIN A2.SS1.p3
 
\begin{center}\includegraphics[width=0.98\linewidth,height=0.72\textheight,keepaspectratio]{rules-09.pdf}\end{center}
 
% END A2.SS1.p3
 
% END A2.SS1
 
% BEGIN A2.SS2
 
\par\medskip\noindent\textbf{B.2 Untyped algorithmic conversion}\par
 
% BEGIN A2.SS2.p1
 
\begin{center}\includegraphics[width=0.98\linewidth,height=0.72\textheight,keepaspectratio]{rules-10.pdf}\end{center}
 
% END A2.SS2.p1
 
% BEGIN A2.SS2.p2
 
\par
 
% END A2.SS2.p2
 
% END A2.SS2
 
% END A2
